\section{The machine model and elementary streams}
\label{sec:streams}

A machine has one fixed finite alphabet and one fixed finite number of
one-dimensional tapes. A step reads and writes at most one cell on each
tape and moves each head by at most one cell. Fixed tracks and finitely
many work tapes change time by a fixed factor. We distinguish moving one
digit in a fixed radix, which needs a bounded number of scans, from
interchanging fields whose numbers of digits grow. The latter permutation
will be realized in Section~\ref{sec:tape-interchange} by computing XOR
combinations of intermediate data streams. The present section supplies
the elementary scans it uses. Every scan includes its counters and the
return of its tape heads.

\subsection{Array order and elementary scans}

For a positive integer $a$ write $[a]=\{0,\ldots,a-1\}$.
An array with address set $[a_1]\times\cdots\times[a_t]$ is stored in
lexicographic order, with the first coordinate most significant.  Its
entries are bits unless a record width is specified.  A trailing record
of $B$ bits can equivalently be regarded as a final address field $[B]$.
An address is implicit in an entry's position in this ordered stream; it
is not stored beside every data bit.
The \emph{logical volume} $V$ is the number of stored data bits.  Work
tapes and duplicate streams are charged in time but do not increase $V$
in recurrences.  All the fields discussed below are nonempty.  A call
with an empty array is handled directly.

An address permutation $f$ moves the entry at $x$ to $f(x)$.  Thus its
action on an array $A$ is $(fA)(f(x))=A(x)$.  Adjacent fields which are
irrelevant to an operation may be replaced by a single field whose length
is their product.  In particular, the elementary procedures below have a
fixed number of formal address fields even when a surrounding algorithm
has a growing list of coordinates.  Lengths and positions are supplied in
canonical binary on descriptor tapes; no additional tape head is attached to an
address coordinate. Each elementary call receives a fixed number of
collapsed canonical lengths. Constructing them from a longer surrounding
list is work of the caller; recursive calls below construct their own
collapsed descriptors within their stated bounds.

When one address field is designated as a row field, collapse the fields
before it to a prefix $[P]$ and those after it to a suffix $[B]$.  If the
row field has range $[R]$, a \emph{row} at prefix $p$ and row index $r$
is the complete consecutive block $\{p\}\times\{r\}\times[B]$.
Every row therefore contains $B$ bits in the same suffix order.

For headers that concatenate a list of positive integer parameters, we use
the self-delimiting code
\begin{equation}
 b_v=\lfloor\log_2v\rfloor+1,\qquad
 \Gamma(v)=0^{b_v-1}\operatorname{bin}_{b_v}(v),\qquad
 |\Gamma(v)|=2b_v-1\quad(v\ge1).
 \label{eq:integer-descriptor-code}
\end{equation}
The initial zeroes specify how many further bits follow the first one,
so concatenated codes can be read in order. We use this code when an
interface specifies a header; elementary calls still receive the fixed
number of canonical binary lengths described above.

\begin{lemma}[Elementary stream operations]
\label{lem:elementary-streams}
Fix a radix $q\ge2$.  Each of the following operations on an array of
volume $V$ takes $O(V)$ steps on a fixed number of one-dimensional tapes:
\begin{enumerate}
\item any fixed permutation of the pair of digit values at two specified
radix-$q$ positions, whether or not the positions are adjacent;
\item moving one specified radix-$q$ digit to another position, preserving
the order of the intervening positions;
\item splitting on a specified radix-$q$ digit, applying a fixed
pointwise map to aligned streams, and merging them back;
\item for positive integers $P,R,B$, within each prefix of an array on
$[P]\times[R]\times[B]$, splitting
whole rows cyclically into a fixed number $W$ of streams, and, when $W$
divides $R$, merging those streams in any fixed order of the row classes.
\end{enumerate}
The bounds include shape preparation, counter initialization and resets,
and returning or clearing the visited portions of work tapes.  Their
constants may depend on the fixed radix, the fixed pointwise map, and the
fixed number of row classes.
\end{lemma}

\begin{proof}
For the first operation split into $q^2$ streams according to the two
specified digits.  With these digits fixed, the order of all remaining
digits is unchanged by the permutation.  Enumerate output addresses in
their required order and read the stream prescribed by the two preimage
digits.  For a digit move use $q$ streams instead.  The same observation
shows that splitting on one digit gives streams with the same remaining
address set in the same order.  Scan their aligned entries, apply the
fixed pointwise map, and merge according to the output digit; this proves
the third operation.  For the fourth, write $r=Wg+w$, where
$0\le w<W$, and send the whole row to stream $w$.  If $W$ divides $R$,
each stream has address set $[P]\times[R/W]\times[B]$, indexed by the
prefix, $g$, and the unchanged suffix.  Thus the streams have the same
shape and each has volume $V/W$.  The inverse merge, with any fixed
permutation of the classes, reads them in the specified cyclic order.

Here are the counter details used in these scans and later ones.  A
binary counter is stored with its head at the least significant end.
For a counter of width $w$, incrementing or decrementing and returning
this head over $L$ consecutive counts costs $O(L+w)$: the number of
visits to position $j$ is $O(1+L/2^j)$. The width term is paid when
initializing the counter and is absorbed by its complete range scan.  A countdown of
length $L$ can be initialized, reset, and tested for underflow in
$O(\log(2L))=O(L)$ time.  Keep a local copy of its length on a dedicated
track, so a reset does not search a larger descriptor.  For a nested scan
of a fixed number of nonempty ranges, the reset of a range is charged to
the scan of that range.  Induction on the number of ranges therefore gives
linear total counting cost.  A fixed number of empty pieces in a split
can be charged to the adjacent nonempty combined range.

Only a fixed number of collapsed lengths are needed.  Computing these
lengths and the initial counter values by ordinary binary arithmetic
takes $O((\log(2V))^C)$ steps for a fixed $C$.  This is $O(V)$ uniformly
in $V\ge1$, after increasing a fixed constant.  Each stream is scanned a
fixed number of times, and every rewind, copy, or erasure is over a
region whose length was charged to that scan.  Delimiters and the fixed
number of tracks require only a fixed alphabet enlargement.
\end{proof}

The next operations are linear because their target field follows their
control fields.  This order will be preserved whenever they are invoked.

\begin{lemma}[Controlled shifts and fixed rational scalings]
\label{lem:ordered-affine-streams}
Fix a prime $q$, a fixed number of control fields, and finitely many
rational coefficients whose denominators are prime to $q$.  Whenever a
coefficient is used as a unit, also require its numerator to be prime to
$q$.  Put $Q=q^b$, $b\ge1$.
Suppose a target field $y\in[Q]$ follows its control fields, each control
used in an offset having $O(b)$ digits.  An update
\[
 y\longmapsto y+a(\text{controls})\pmod Q
\]
costs $O(V)$ if the offset can be prepared in $b^{O(1)}$ time from the
fixed number of control values.  In particular, this holds for a fixed
rational linear combination of them, specialized modulo $Q$.
Multiplication of $y$ by any of the specified rational units, and its
inverse, also costs $O(V)$.  Arbitrarily long intervening spectator
fields and suffix records are permitted.
\end{lemma}

\begin{proof}
Fix all fields before the target, and let the suffix length be $B$ bits.
The resulting fiber consists of $Q$ consecutive blocks of length $B$.
For an offset $a\in[Q]$, split these blocks at $Q-a$.  Append the two
pieces to two tapes, doing so successively for all prefix fibers.  In the
merge read the second piece and then the first piece for each prefix.
The pieces have become the two consecutive ranges $[0,a)$ and $[a,Q)$
of the new target coordinate.  Recompute $a$ when starting that prefix in
the merge.  Initializing these two piece lengths and returning the heads
of the relevant control counters is included in the $b^{O(1)}$ offset
work.  Since
\[
 b^C=O(q^b)=O(Q)
\]
for each fixed $C$, this work is absorbed by the $QB\ge Q$ bits of the
fiber.  Unrelated spectator counters are incremented by ripple counting;
their entire contents are not consulted when computing the offset.

For positive fixed integer $c$ prime to $q$, the permutation
$y\mapsto cy\pmod Q$ has $c$ increasing pieces, classified by
\[
 j=\left\lfloor\frac{cy}{Q}\right\rfloor,
 \qquad 0\le j<c.
\]
Prepare their endpoints once for this $Q$, split each fiber into these
$c$ streams, and enumerate the output coordinate $y'=0,\ldots,Q-1$.
The required stream is the unique $j\in[c]$ for which
$y'+jQ\equiv0\pmod c$.  Uniqueness follows from $\gcd(Q,c)=1$;
the corresponding $y=(y'+jQ)/c$ lies in $[Q]$.  The stream choice is
maintained by a finite residue counter.  Reversing the roles of this
split and merge implements multiplication by $c^{-1}$ modulo $Q$.

Negation leaves the block at zero fixed and reverses the order of the
other $Q-1$ blocks.  Push these blocks consecutively to a temporary tape.
Popping gives their order reversed, but also reverses the bits in each
block.  Reverse each popped block on a second temporary tape before
writing it, restoring its internal order.  Every bit is moved a bounded
number of times.  The same procedure is repeated for successive fibers;
the temporary stack is emptied and its head returned locally.  A signed
rational unit is a composition of these integer multiplications, inverse
multiplications, and possibly negation.  Its description is fixed.

The nested scan and reset analysis of
Lemma~\ref{lem:elementary-streams} applies to all these operations.  In
particular, copying a suffix-length counter at every target position
costs $O(\log(2B))\le O(B)$ at that position.  It does not introduce a
factor $\log V$ per bit.
\end{proof}

When coefficients have a denominator $d$ prime to $q$, their reductions
modulo $q^b$ can be computed by the extended Euclidean algorithm using
ordinary tape arithmetic.  The dimension and all original coefficients
are fixed, so these preparations have polynomial cost in $b$.  Nothing
in Lemma~\ref{lem:ordered-affine-streams} assumes unit-cost arithmetic.

For later use, fixed-width signed fixed-point binary addition,
subtraction, sign change, and multiplication or division by a fixed
power of two are streaming operations of linear bit cost, provided
enough integer and fractional guard bits have been allocated for exact
results. If an input record is most-significant-bit first, reverse it onto
the arithmetic work tape. With least significant bits first on that tape,
propagate carries using finite control. Reverse the completed record if
its required external order is opposite.
This convention separates the order of bits inside a record from the
lexicographic order of records. Each required reversal is a linear pass
and is charged.
